\documentclass[UTF8,heading=true,12pt,font=songti]{article}
\usepackage{ctex}
\usepackage{cite}
\usepackage{fontawesome}
\usepackage{pdfpages}
\usepackage{titlesec}
\usepackage{xparse,amsmath,amsfonts,mathrsfs,amsthm,amssymb,tikz,bm,graphicx,url,listings,subfig,float,subcaption}
\usepackage{physics}
%\usepackage{mathptmx}% 新罗马字体
\usepackage[ruled,linesnumbered]{algorithm2e}
\usepackage[colorlinks,linkcolor=blue]{hyperref} % black 
\usepackage{gbt7714}% 参考文献格式
\usepackage[left=3cm,right=2.5cm,top=2.5cm,bottom=2.5cm]{geometry}% ,showframe 显示框图
\usepackage{lipsum}% 生成示例
\usepackage{titlesec,titletoc} % 设置目录格式
\renewcommand{\contentsname}{\centering{\quad \\
		\heiti\zihao{3}\bfseries{目\qquad 录\par\par}}} % 黑体三号字加黑居中
\usepackage[titles,subfigure]{tocloft}
\renewcommand\cftsecdotsep{\cftdotsep}
\renewcommand\cftsecleader{\cftdotfill{\cftsecdotsep}}
\renewcommand\cftsecfont{\songti\zihao{-4}} % section字号字体
\renewcommand\cftsubsecfont{\songti\zihao{-4}} % subsection字号字体
\renewcommand\cftsubsubsecfont{\songti\zihao{-4}} % subsubsection字号字体
\renewcommand\cftsecpagefont{\mdseries}
\renewcommand\cftsubsecpagefont{\mdseries}
\renewcommand{\bibfont}{\zihao{5}}

%\titlecontents{section}[1.5em]{\zihao{-4}\bfseries}{\contentslabel[\thecontentslabel]{1.5em}}{}{\mdseries\titlerule*[2ex]{.}\contentspage}
%定义环境
\renewcommand{\qed}{\hfill\blacksquare}
\newcommand{\qedwhite}{\hfill \ensuremath{\Box}}
\numberwithin{equation}{subsection}
\newtheorem{theorem}{\hspace{2em}定理}[subsection]
\newtheorem{lemma}{\hspace{2em}引理}[subsection]
\newtheorem{definition}{\hspace{2em}定义}[subsection]
\newcommand{\pf}{\textbf{证明}$\quad$}
\newcommand{\q}{\textbf{问题}$\quad$}
\newcommand{\s}{\textbf{解答}$\quad$}

\usepackage{setspace}
\linespread{1.5} % 全局1.5倍行间距
\pagestyle{plain} % 仅页脚显示页数
\everymath{\displaystyle}
% 代码格式
\lstset{
	frame=trbl,
	aboveskip=3mm,
	belowskip=3mm,
	showstringspaces=true,
	breaklines=true,
	columns=flexible,
	framerule=1pt,
	rulecolor=\color{gray},
	backgroundcolor=\color{white},
	basicstyle={\small\ttfamily},
	numbers=left,
	numbersep=1em,
	numberstyle=\color{black},
	keywordstyle=\color{blue},
	commentstyle=\color{gray},
	stringstyle=\color{mauve},
	tabsize=3,
}

% 节标题字体设置
\ctexset{section={format={\songti \zihao{4}\bfseries}},
	subsection={format={\songti \zihao{-4}\bfseries}},
	subsubsection={format={\fangsong \zihao{-4}}}}

% 表格及图片标题样式设置-
\renewcommand{\thefigure}{\arabic{section}-\arabic{figure}} %使图编号为1-1的格式 %\protect{~}
\renewcommand{\thetable}{\arabic{section}-\arabic{table}}%使表编号为1-1的格式

%表格与图片标题字体设置宋体小五
\captionsetup[figure]{textfont=footnotesize,labelfont=footnotesize,labelformat=simple, labelsep=space}
\captionsetup[subfigure]{labelfont=footnotesize,textfont=footnotesize}
\captionsetup[table]{labelfont=footnotesize,textfont=footnotesize,labelformat=simple, labelsep=space} % 论文格式设置文件

%\title{\songti \zihao{3}\bfseries 标题}
%\author{}
%\date{\today}

\begin{document}
    %\maketitle% 标题页
    \includepdf[pages=-]{fig/titlepage}
    \newpage
    \tableofcontents% 生成目录
    \setcounter{page}{0}% 此为第零页
    \thispagestyle{empty}% 页眉页脚为空
    \newpage
    \begin{center}
    \songti \zihao{3}\bfseries 标题
	\end{center}

    \begin{table}[H]
        \vspace{-1.5em}
        \centering
        \begin{spacing}{1.5}
        \songti \zihao{5}
        \begin{tabular}{rl} %设置表格字体位置(右左)
            学生姓名:\ NAME & 学号:\ 20215034000\\ 
            数学与统计学院 & 信息与计算科学专业\\
            指导教师:\ NAME & 职称:\ \\	
        \end{tabular}
        \end{spacing}
        \vspace{-1.5em}
    \end{table}
	
    \addcontentsline{toc}{section}{摘\quad 要}\tolerance=500 % 将摘要放进目录
    \noindent {\heiti \zihao{-4}\bfseries{摘\quad 要}}: \kaishu\zihao{-4} \ 摘要内容 % 无缩进


    \noindent {\heiti \zihao{-4}\bfseries{关键词}}:\kaishu\zihao{-4} \ 内容 % 无缩进
    
    \begin{center}
    	\zihao{3} Title
    \end{center}
    \addcontentsline{toc}{section}{\bfseries Abstract}\tolerance=500 % 将Abstract放进目录
    \noindent {\zihao{-4}\bfseries{Abstract}}:\ \zihao{-4} \lipsum[1]
    
    \noindent {\zihao{-4}\bfseries{Keywords}}:\ \zihao{-4} Ab
    
    %\onehalfspacing %1.5倍行间距
    \section{前言}
    \songti
	$$
    f(z_0)=\frac{1}{2\pi i}\oint_{\Gamma}\frac{f(z)}{z-z_0}\dd z.
    $$

    \section{第一节}
    \lipsum[3]
    \begin{table}[H]
        \centering
        \caption{符号说明}
        \begin{tabular}{clc} %设置表格字体位置(中左中)
            \hline
            \textbf{符号} & \textbf{说明} & \textbf{单位} \\ \hline
            $\eta$ & 定日镜的光学效率 & 无量纲\\	
            \hline
        \end{tabular}
    \end{table}

    \subsection{第一小节}
    引入参考文献\cite{bazlov2024momentsrepresentationnumbers}


    \subsubsection{小小一节}
    \begin{figure}[H]
		\centering
		\subfloat[示意图]{\includegraphics[width = 0.5\linewidth]{figure/logo-r.eps}}
		\hfill
		\subfloat[示意图]{\includegraphics[width = 0.5\linewidth]{figure/logo-r.eps}}
		\hfill
		\subfloat[示意图]{\includegraphics[width = 0.5\linewidth]{figure/logo-r.eps}}
		\hfill
		\subfloat[示意图]{\includegraphics[width = 0.5\linewidth]{figure/logo-r.eps}}
		\caption{示意图}
		\label{fig:label2}
	\end{figure}  

    \bibliography{ref/ref}
    \addcontentsline{toc}{section}{参考文献}\tolerance=500 % 将参考文献放进目录

\end{document}
    
